\documentclass[12pt,a4paper,openany]{report}

% ── Packages ──────────────────────────────────────────────────────────────────
\usepackage[utf8]{inputenc}
\usepackage[T1]{fontenc}
\usepackage{lmodern}
\usepackage[margin=2.5cm]{geometry}
\usepackage{setspace}
\usepackage{parskip}
\usepackage{enumitem}
\usepackage{booktabs}
\usepackage{longtable}
\usepackage{array}
\usepackage{tabularx}
\usepackage{multirow}
\usepackage{amsmath,amssymb}
\usepackage{xcolor}
\usepackage{listings}
\usepackage{fancyvrb}
\usepackage{float}
\usepackage{caption}
\usepackage{subcaption}
\usepackage{titlesec}
\usepackage{titletoc}
\usepackage{hyperref}
\usepackage{etoolbox}
\usepackage{microtype}
\usepackage{verbatim}
\usepackage{alltt}

% ── Hyperref ──────────────────────────────────────────────────────────────────
\hypersetup{
  colorlinks=true,
  linkcolor=blue!70!black,
  citecolor=blue!70!black,
  urlcolor=blue!70!black,
  pdftitle={Retainly: Offline-First Adaptive Study Planner — Technical White Paper},
  pdfauthor={CodeSym},
  pdfsubject={Technical White Paper},
  pdfkeywords={Flutter, offline-first, study planner, SM-2, Firebase, SQLite}
}

% ── Listing style ─────────────────────────────────────────────────────────────
\lstset{
  basicstyle=\small\ttfamily,
  breaklines=true,
  frame=single,
  rulecolor=\color{gray!30},
  columns=fullflexible,
  keepspaces=true
}

% ── Verbatim style ────────────────────────────────────────────────────────────
\DefineVerbatimEnvironment{blockverbatim}{Verbatim}{fontsize=\small,frame=single}

% ── Section formatting ────────────────────────────────────────────────────────
\titleformat{\chapter}[display]
  {\normalfont\Huge\bfseries}
  {\chaptertitlename\ \thechapter}{20pt}{\Huge}
\titlespacing*{\chapter}{0pt}{-30pt}{40pt}

\titleformat{\section}
  {\normalfont\Large\bfseries}
  {\thesection}{1em}{}
\titleformat{\subsection}
  {\normalfont\large\bfseries}
  {\thesubsection}{1em}{}
\titleformat{\subsubsection}
  {\normalfont\normalsize\bfseries}
  {\thesubsubsection}{1em}{}

% ── Title page ────────────────────────────────────────────────────────────────
\title{
  \textbf{\Huge Retainly}\\[0.5em]
  \Large Offline-First Adaptive Study Planner\\[2em]
  \large Technical White Paper
}
\author{
  CodeSym\\
  \texttt{https://github.com/CodeSym/retainly}
}
\date{Version 1.0 --- \today}

% ── Begin document ────────────────────────────────────────────────────────────
\begin{document}

\maketitle
\thispagestyle{empty}
\newpage

\begin{abstract}
\noindent
Retainly is a cross-platform, offline-first mobile application engineered to deliver adaptive daily study planning, Pomodoro-style focus tracking, and spaced-repetition revision for Pakistani Matriculation (Class 9--10) students. The system is architected around a local SQLite single source of truth with optional Firebase-backed cloud synchronization, AI-assisted task decomposition, and encrypted local backup. This white paper exhaustively documents the system's layered architecture, data-flow invariants, algorithmic mechanics, concurrency model, security posture, and extensibility roadmap.
\end{abstract}

\newpage
\tableofcontents
\newpage

% ═══════════════════════════════════════════════════════════════════════════════
\chapter{Executive Summary \& Core Objectives}
% ═══════════════════════════════════════════════════════════════════════════════

\section{Mission Statement}

Retainly exists to transform a static syllabus into an adaptive, personalized daily study plan for students in low-connectivity environments. The system must function identically on a device with no network access as it does online, with zero data loss and no crash paths attributable to unavailable cloud services.

\section{Core Technical Goals}

\begin{table}[H]
\centering
\caption{Core technical goals and their implementation mechanisms}
\label{tab:core-goals}
\begin{tabular}{>{\bfseries}l p{10cm}}
\toprule
\textbf{Goal} & \textbf{Implementation Mechanism} \\
\midrule
Offline-first resilience & Local SQLite as sole source of truth; all Firebase dependencies degrade silently \\
Adaptive planning & Greedy knapsack daily budget allocation scored by urgency, priority, and fit \\
Retention optimization & SuperMemo-2 (SM-2) spaced-repetition algorithm with adaptive ease factor \\
Data durability & Encrypted AES-256 local backups; sync outbox with tombstone-based conflict prevention \\
Cross-platform parity & Flutter 3.24+ targeting Android (primary) and Linux desktop; explicit web rejection \\
Privacy-by-default & AI/OCR/sync strictly opt-in with consent, cost-warning, and daily-quota gates \\
Academic integrity & AI outputs labeled as draft-only; hallucination reporting; source-grounding prompts \\
\bottomrule
\end{tabular}
\end{table}

\section{Non-Goals}

\begin{itemize}[leftmargin=*]
  \item Multi-user collaboration or real-time co-editing
  \item Web platform support (explicitly rejected; \texttt{firebase\_options.dart} throws \texttt{UnsupportedError} on web)
  \item Native iOS or macOS as primary targets (desktop FFI fallback only)
  \item Offline semantic merge of conflicting remote mutations (manual resolution required)
\end{itemize}

\section{Architectural Philosophy}

The system follows a \textbf{layered, feature-first, domain-driven} architecture with three immutable invariants:

\begin{enumerate}[leftmargin=*]
  \item \textbf{Local-first authority:} The SQLite database is the single source of truth. No UI state is derived directly from cloud state.
  \item \textbf{No Firebase crash path:} Every Firebase-dependent service catches \texttt{FirebaseException} and returns a safe default. The application never terminates due to cloud unavailability.
  \item \textbf{Explicit opt-in for online features:} Cloud sync, AI assistance, OCR, analytics, and crash reporting require explicit user consent and fail gracefully when denied or unavailable.
\end{enumerate}

\section{System Invariants}

\begin{verbatim}
INV-1: Every local mutation flows through DatabaseRepository.
INV-2: Every local mutation is recorded in sync_meta via _trackChange().
INV-3: No UI screen depends on Firebase availability for rendering.
INV-4: AI/OCR requests are gated by (consent AND cost_warning AND quota AND connectivity).
INV-5: Focus session state transitions are: idle -> running -> completed|idle.
INV-6: Sync tombstones are retained >=30 days to prevent resurrection attacks.
INV-7: SM-2 ease factor is clamped to [1.3, infinity).
INV-8: Backup encryption key is stored in platform secure storage when available.
\end{verbatim}

% ═══════════════════════════════════════════════════════════════════════════════
\chapter{System Architecture \& High-Level Design}
% ═══════════════════════════════════════════════════════════════════════════════

\section{Architectural Paradigm}

Retainly employs a \textbf{layered, offline-first, event-driven} architecture. The client is a single Flutter application with a clear separation between presentation, service, and data layers. The backend is a thin Firebase layer consisting of four Cloud Functions and Firestore/Storage security rules.

\begin{figure}[H]
\centering
\begin{blockverbatim}
+-------------------------------------------------------------+
|                      PRESENTATION LAYER                     |
|  Flutter Widgets (Material 3) + Riverpod State + GoRouter   |
|  +----------+ +----------+ +----------+ +----------+ +------+ |
|  |   Home   | | Planner  | |  Focus   | | Revision | |Backup| |
|  | Screen   | | Screen   | | Screen   | | Screen   | |Screen| |
|  +----+-----+ +----+-----+ +----+-----+ +----+-----+ +--+---+ |
|       |            |            |            |            |    |
|  +----+------------+------------+------------+------------+---+ |
|  |              Riverpod Providers (DI + State)             | |
|  | dbFutureProvider | userProfileProvider | dashboardProvider | |
|  +------------------------+-----------------------------------+ |
+---------------------------+-----------------------------------+
                            |
+---------------------------+-----------------------------------+
|                        SERVICE LAYER                         |
|  +-------------+ +-------------+ +-------------+ +-----------+ |
|  | SyncService | |OfflineQueue | |  AIService  | |Connectivity| |
|  |(singleton)  | |  Service    | |(singleton)  | | Service    | |
|  +------+------+ +------+------+ +------+------+ +------+----+ |
|         |             |             |             |           |
|  +------+-------------+-------------+-------------+-----------+ |
|  |              SyncWorkerService (Workmanager)              | |
|  |        Background Isolate @pragma('vm:entry-point')       | |
|  +------------------------+----------------------------------+ |
+---------------------------+-----------------------------------+
                            |
+---------------------------+-----------------------------------+
|                         DATA LAYER                            |
|  +----------------------------------------------------------+  |
|  |                  DatabaseRepository (Facade)             |  |
|  | Wraps: DatabaseHelper (sqflite) + AppDatabase (Drift)    |  |
|  | Responsibilities: CRUD, SM-2, analytics, sync tracking   |  |
|  +------------------------+---------------------------------+  |
|                            |                                  |
|  +------------------------+---------------------------------+  |
|  |                DatabaseHelper (sqflite)                  |  |
|  |   11-table schema, singleton, backward-compatible mig.   |  |
|  |   File: getApplicationDocumentsDirectory()/study_planner.db| |
|  +----------------------------------------------------------+  |
|                                                                 |
|  +----------------------------------------------------------+  |
|  |             AppDatabase (Drift, Android-only)            |  |
|  |   schemaVersion: 3, NativeDatabase.createInBackground()  |  |
|  +----------------------------------------------------------+  |
+-------------------------------------------------------------+
                            |
                +-----------+-----------+
                |   Local Only (Default) |
                |   OR                   |
                |   Firebase (Optional)  |
                +-----------+-----------+
                            |
+---------------------------+-----------------------------------+
|                    CLOUD LAYER (Optional)                      |
|  +----------------------------------------------------------+  |
|  |                  Firebase Firestore                       |  |
|  | Collections: sync_outbox, sync_tombstones, sync_conflicts,| |
|  | ai_consents, ai_requests, ocr_jobs, quizzes, quiz_attempts| |
|  +------------------------+---------------------------------+  |
|                            |                                   |
|  +------------------------+---------------------------------+  |
|  |         Firebase Cloud Functions (Node.js/TS)            |  |
|  |  +---------+ +---------+ +-------------+ +------------+  |  |
|  |  | aiProxy | |ocrProcess| | syncWorker  | |pruneTombstones| |
|  |  | (https) | | (https) | |(pubsub 15m) | | (pubsub 24h)|  |  |
|  |  +---------+ +---------+ +-------------+ +------------+  |  |
|  +----------------------------------------------------------+  |
|                                                                 |
|  +----------------------------------------------------------+  |
|  |           Firebase Storage (ai-uploads/, ocr-output/)    |  |
|  +----------------------------------------------------------+  |
+-------------------------------------------------------------+
\end{blockverbatim}
\caption{Retainly full-system layered architecture}
\label{fig:system-architecture}
\end{figure}

\section{Module Boundaries \& Dependencies}

\begin{blockverbatim}
lib/
+-- core/                  # Zero dependencies on data/services
|   +-- constants/         # AppConstants, MatricSubjects
|   +-- enums/             # ResourceType
|   +-- feature_flags.dart # Compile-time/runtime feature toggles
|   +-- theme/             # AppTheme (Material 3)
|   +-- utils/             # error_utils, planner_utils
|
+-- data/                  # Depends only on core
|   +-- models/            # Plain Dart data classes (toMap/fromMap)
|   +-- database_helper.dart    # sqflite singleton
|   +-- drift/             # Optional Drift ORM (Android only)
|   +-- firestore/         # Firestore-specific models
|   +-- repositories/      # DatabaseRepository (facade)
|
+-- services/              # Depends on data + Firebase SDKs
|   +-- sync_service.dart
|   +-- offline_queue_service.dart
|   +-- sync_worker_service.dart
|   +-- ai_service.dart
|   +-- connectivity_service.dart
|   +-- analytics_service.dart
|   +-- crashlytics_service.dart
|   +-- remote_config_service.dart
|   +-- shortcut_service.dart
|
+-- features/              # Depends on services + data + providers
|   +-- ai/                # OCR, flashcards, quiz, hallucination reports
|   +-- backup/            # Full DDD: domain/data/application/presentation
|   +-- focus/             # Pomodoro timer + Android DND channel
|   +-- planner/           # Daily plan generation + reschedule
|   +-- revision/          # SM-2 spaced repetition queue
|   +-- resources/         # PDF viewer, practical records
|   +-- subjects/          # Subject setup, chapter management
|   +-- tasks/             # Add task, global search
|   +-- ...                # home, progress, settings, legal, onboarding
|
+-- navigation/            # GoRouter config (depends on providers)
+-- providers/             # Riverpod providers (DI + state)
+-- l10n/                  # Generated localization (en + ur)
\end{blockverbatim}

\section{Inter-Process Communication}

\begin{table}[H]
\centering
\caption{Inter-process communication mechanisms}
\label{tab:ipc}
\begin{tabular}{llp{7cm}}
\toprule
\textbf{Mechanism} & \textbf{Direction} & \textbf{Purpose} \\
\midrule
\texttt{MethodChannel('focus\_shield')} & Dart $\leftrightarrow$ Kotlin & Toggle Android DND mode during focus sessions \\
\texttt{MethodChannel('app\_shortcuts')} & Dart $\leftrightarrow$ Kotlin & Handle Android app shortcut launch intent \\
Workmanager background isolate & Main isolate $\leftrightarrow$ Background isolate & Periodic sync task (15 min) in separate Dart VM \\
\texttt{StreamController.broadcast()} & Service $\rightarrow$ UI & Offline queue change notifications \\
\texttt{connectivity\_plus} stream & Native $\rightarrow$ Dart & Connectivity state changes (online/offline banner) \\
SharedPreferences & Main isolate $\leftrightarrow$ Background isolate & Persistent config (quota counters, queue state) \\
\bottomrule
\end{tabular}
\end{table}

\section{Onboarding State Machine}

The navigation layer implements a two-stage gate implemented via \texttt{\_RouterRefreshNotifier}:

\begin{figure}[H]
\centering
\begin{blockverbatim}
+-------------+
|   Start     |
+------+------+
       |
       v
+---------------------+
| dbFutureProvider?   |-- No --> return null (wait)
|   resolved?         |
+----------+----------+
           | Yes
           v
+---------------------+
| userProfile exists? |-- No --> redirect /onboarding
|                     |
+----------+----------+
           | Yes
           v
+---------------------+
| Normal Navigation   |
| (ShellRoute / modal)|
+---------------------+
\end{blockverbatim}
\caption{Onboarding state machine}
\label{fig:onboarding-state-machine}
\end{figure}

\section{Focus Session State Machine}

The \texttt{FocusScreen} implements a focus-session state machine with three canonical states:

\begin{figure}[H]
\centering
\begin{blockverbatim}
+---------+
|  IDLE   | <- New session can be initiated
+----+----+
     | startSession()
     v
+---------+
| RUNNING | <- Timer active, DND enabled (Android)
+----+----+
     | complete() / pause() / app dispose()
     v
+---------+
|COMPLETED| <- Session persisted, DND disabled
+----+----+
     | reset()
     v
+---------+
|  IDLE   |
+---------+
\end{blockverbatim}
\caption{Focus session lifecycle state machine}
\label{fig:focus-state-machine}
\end{figure}

On \texttt{dispose()}, if a session is still running, \texttt{\_saveSession()} persists the in-progress session to prevent data loss on app exit. The Android platform channel (\texttt{focus\_shield}) toggles Do Not Disturb via \texttt{NotificationManager.setInterruptionFilter(INTERRUPTION\_FILTER\_NONE)}, guarded by a \texttt{FOREGROUND\_SERVICE} (\texttt{FocusShieldForegroundService}).

% ═══════════════════════════════════════════════════════════════════════════════
\chapter{Deep-Dive Subsystem Mechanics}
% ═══════════════════════════════════════════════════════════════════════════════

\section{Data Layer}

\subsection{DatabaseHelper (sqflite Singleton)}

\texttt{DatabaseHelper} implements a thread-safe singleton with double-checked locking:

\begin{lstlisting}[language=Java,caption={DatabaseHelper double-checked locking pattern}]
class DatabaseHelper {
  static final DatabaseHelper instance = DatabaseHelper._init();
  static Database? _database;
  static Future<Database>? _databaseOpening;

  DatabaseHelper._init();

  Future<Database> get database async {
    if (_database != null) return _database!;
    _databaseOpening ??= _initDB('study_planner.db');
    _database = await _databaseOpening!;
    return _database!;
  }
}
\end{lstlisting}

The database file resides at \texttt{getApplicationDocumentsDirectory()/study\_planner.db}. On creation, 11 tables are defined with foreign-key constraints (\texttt{ON DELETE CASCADE} / \texttt{ON DELETE SET NULL}).

\paragraph{Schema versioning strategy.}
The database is opened at \texttt{version: 1}. On every open, \texttt{\_extendSchema()} runs inside a transaction, querying \texttt{PRAGMA table\_info()} for each table to detect missing columns and applying \texttt{ALTER TABLE} additions idempotently. This enables backward-compatible schema evolution without data destruction.

\paragraph{11-Table Schema.}

\begin{table}[H]
\centering
\caption{SQLite schema overview}
\label{tab:schema}
\small
\begin{tabular}{llll}
\toprule
\textbf{Table} & \textbf{Primary Key} & \textbf{Foreign Keys} & \textbf{Purpose} \\
\midrule
\texttt{user\_profiles} & \texttt{id} (auto) & --- & Single-row user configuration \\
\texttt{subjects} & \texttt{id} (auto) & --- & Academic subjects \\
\texttt{chapters} & \texttt{id} (auto) & \texttt{subject\_id} $\to$ \texttt{subjects} & Chapter-level granularity \\
\texttt{study\_tasks} & \texttt{id} (auto) & \texttt{subject\_id}, \texttt{chapter\_id} & Scheduled study tasks \\
\texttt{focus\_sessions} & \texttt{id} (auto) & \texttt{task\_id} $\to$ \texttt{study\_tasks} & Pomodoro session records \\
\texttt{revision\_items} & \texttt{id} (auto) & \texttt{chapter\_id} $\to$ \texttt{chapters} & Spaced-repetition queue \\
\texttt{resources} & \texttt{id} (auto) & \texttt{subject\_id}, \texttt{chapter\_id} & PDF/pinned resources \\
\texttt{practical\_records} & \texttt{id} (auto) & \texttt{subject\_id}, \texttt{resource\_id} & Lab practical records \\
\texttt{backup\_records} & \texttt{id} (auto) & --- & Backup audit trail \\
\texttt{syllabus\_templates} & \texttt{id} (auto) & --- & Imported/exported syllabus schemas \\
\texttt{sync\_meta} & \texttt{id} (auto) & --- & Local outbox metadata; \texttt{UNIQUE(entity, local\_id)} \\
\bottomrule
\end{tabular}
\end{table}

\subsection{Drift ORM (Optional, Android-Only)}

\texttt{AppDatabase} provides type-safe DAO access via the Drift framework. It is initialized only when \texttt{Platform.isAndroid} using \texttt{NativeDatabase.createInBackground()}, which offloads SQLite I/O to a background isolate. The \texttt{DatabaseRepository} selects between sqflite and Drift at runtime via the \texttt{\_useDrift} getter:

\begin{lstlisting}[language=Java,caption={Runtime Drift/sqflite selection in DatabaseRepository}]
bool get _useDrift => driftDb != null;

Future<int> insertSubject(SubjectModel subject) async {
  int id;
  if (_useDrift) {
    final db = driftDb!;
    id = await db.into(db.subjects).insert(
      SubjectsCompanion.insert(
        name: subject.name,
        color: subject.color,
        sortOrder: subject.sortOrder,
        createdAt: subject.createdAt,
      ),
    );
  } else {
    id = await db.insertSubject(subject.toMap());
  }
  await _trackChange('subjects', id.toString(), 'create', subject.toMap());
  return id;
}
\end{lstlisting}

The Drift schema is maintained in parallel with the raw SQLite schema; both must be updated when columns or tables change. \texttt{schemaVersion} is currently 3.

\subsection{Data Models}

All models are plain Dart classes with \texttt{toMap()} and \texttt{fromMap()} constructors. \texttt{fromMap()} uses defensive type checking to handle partial or legacy rows:

\begin{lstlisting}[language=Java,caption={Defensive type checking in UserModel.fromMap()}]
factory UserModel.fromMap(Map<String, dynamic> map) {
  return UserModel(
    id: map['id'] as int?,
    studentName: map['student_name'] is String
        ? map['student_name'] as String
        : '',
    // ... all fields with defensive casting
  );
}
\end{lstlisting}

This protects against:
\begin{itemize}[leftmargin=*]
  \item Rows from older schema versions that lack newer columns
  \item Malformed data from backup imports
  \item Type coercion edge cases between sqflite and Drift
\end{itemize}

\subsection{DatabaseRepository (Central Facade)}

\texttt{DatabaseRepository} is the \textbf{sole intermediary} between the UI layer and the database. It enforces INV-1 and INV-2. Key responsibilities:

\paragraph{CRUD Operations.}
Every entity type has typed accessors (e.g., \texttt{getTodayTasks()}, \texttt{getTasksForDate()}, \texttt{getSubjectProgress()}). Each mutating method calls \texttt{\_trackChange()} to write a \texttt{sync\_meta} entry.

\paragraph{SM-2 Spaced Repetition Algorithm.}

The implementation follows the SuperMemo-2 specification with the following mechanics:

\begin{align}
\text{ease\_factor}_{\text{new}} &= \max\left(1.3,\ \text{ease\_factor}_{\text{current}} + \left(0.1 - (5 - q) \cdot (0.08 + (5 - q) \cdot 0.02)\right)\right) \\
\text{interval}_{\text{new}} &= 
\begin{cases}
1 & \text{if repetitions} = 0 \\
6 & \text{if repetitions} = 1 \\
6 & \text{if repetitions} = 2 \\
\text{clamp}(\text{round}(\text{interval}_{\text{current}} \cdot \text{ease\_factor}_{\text{new}}), 1, 365) & \text{otherwise}
\end{cases}
\end{align}

Confidence-to-quality mapping:
\begin{align}
\text{quality} &= 
\begin{cases}
5 & \text{if confidence} \ge 90 \\
4 & \text{if confidence} \ge 75 \\
3 & \text{if confidence} \ge 60 \\
2 & \text{if confidence} \ge 40 \\
1 & \text{if confidence} \ge 20 \\
0 & \text{otherwise}
\end{cases}
\end{align}

\paragraph{Task Scoring Function.}
\begin{equation}
\text{score} = 
\underbrace{\max(0,\ 30 - \min(30,\ \Delta_{\text{days}}))}_{\text{urgency}} + 
\underbrace{10 \cdot \text{priority}}_{\text{priority}} + 
\underbrace{5 \cdot \mathbb{1}_{\{\text{estimated\_minutes} \le \text{daily\_minutes}\}}}_{\text{fit bonus}}
\end{equation}
where $\Delta_{\text{days}}$ is the non-negative difference between the due date and today (capped at 30), and tasks already overdue receive an additional urgency bonus.

\paragraph{Analytics Queries.}
Six raw SQL analytics methods provide insights:
\begin{itemize}[leftmargin=*]
  \item \texttt{getRecallTrends()} -- 90-day rolling average confidence by subject
  \item \texttt{getSubjectConfidenceDecay()} -- Subject-level confidence aggregation
  \item \texttt{getTaskEstimateAccuracy()} -- Estimated vs.\ actual focus minutes per task
  \item \texttt{getProductiveTimeInsights()} -- Hour-of-day focus session distribution
  \item \texttt{getMissedDayPatterns()} -- 30-day completion rate by date
  \item \texttt{getSubjectEstimateAccuracy()} -- Per-subject estimate accuracy
\end{itemize}

\paragraph{Sync Tracking (\texttt{\_trackChange}).}
\begin{lstlisting}[language=Java,caption={Sync metadata tracking}]
Future<void> _trackChange(String entity, String localId, String operation,
    Map<String, dynamic> data) async {
  final now = DateTime.now().toIso8601String();
  final rawDb = await db.database;
  await rawDb.insert('sync_meta', {
    'entity': entity,
    'local_id': localId,
    'sync_status': 'pending',
    'conflict_data': null,
    'created_at': now,
    'updated_at': now,
  }, conflictAlgorithm: ConflictAlgorithm.replace);
}
\end{lstlisting}

The \texttt{UNIQUE(entity, local\_id)} constraint ensures idempotent retry behavior.

\section{Service Layer}

\subsection{SyncService}

A singleton bridging local changes to Firestore. Key operations:

\begin{itemize}[leftmargin=*]
  \item \texttt{enqueueLocalChange()} -- Writes a \texttt{sync\_outbox} document with \texttt{synced: false}. Does NOT directly mutate the target collection; the Cloud Function \texttt{syncWorker} handles application.
  \item \texttt{markEntityDeleted()} -- Writes a \texttt{sync\_tombstones} document with \texttt{deletedAt} timestamp.
  \item \texttt{retryPending()} -- Client-side fallback that reads unsynced outbox items and applies them directly to target collections. Uses last-write-wins for non-delete operations.
  \item \texttt{pruneTombstones()} -- Deletes tombstones older than 30 days.
  \item \texttt{recordConflict()} -- Writes a \texttt{sync\_conflicts} document with \texttt{localData} and \texttt{remoteData} snapshots.
\end{itemize}

All methods follow the guard pattern:

\begin{lstlisting}[language=Java,caption={Firebase initialization guard pattern}]
Future<void> _ensureFirebase() async {
  if (_firestore != null) return;
  try {
    _firestore = FirebaseFirestore.instance;
  } catch (_) {
    _firestore = null;
  }
}
\end{lstlisting}

And every public method catches \texttt{FirebaseException} silently:
\begin{lstlisting}[language=Java]
try {
  await _firestore!.collection('sync_outbox').add({...});
} on FirebaseException catch (_) {}
\end{lstlisting}

\subsection{OfflineQueueService}

A client-side queue persisted as JSON in \texttt{SharedPreferences} under the key \texttt{offline\_queue}. It holds \texttt{QueuedOperation} objects with fields: \texttt{type} (\texttt{sync}, \texttt{ai}, \texttt{tombstone}), \texttt{payload}, \texttt{attemptCount}, \texttt{createdAt}, \texttt{updatedAt}.

Processing semantics:
\begin{itemize}[leftmargin=*]
  \item \texttt{processQueue()} iterates queued operations, attempting each.
  \item On failure, \texttt{attemptCount} is incremented.
  \item Operations exceeding \texttt{\_maxAttempts} (3) are dropped.
  \item A \texttt{StreamController<void>.broadcast()} notifies listeners.
\end{itemize}

This queue is distinct from the Firestore \texttt{sync\_outbox}: the \texttt{OfflineQueueService} buffers operations when the device is offline; the Firestore outbox is the eventual cloud consistency target.

\subsection{SyncWorkerService}

Orchestrates background sync via \texttt{workmanager}. Registers a periodic task (\texttt{sync\_worker\_process\_queue}) with:
\begin{itemize}[leftmargin=*]
  \item Frequency: 15 minutes
  \item Network constraint: \texttt{NetworkType.connected}
  \item Backoff: linear starting at 5 minutes
\end{itemize}

The worker runs in a \textbf{separate Dart isolate} using \texttt{@pragma('vm:entry-point')} to ensure the entry point is retained in release builds. It initializes Firebase independently from the main isolate:

\begin{lstlisting}[language=Java,caption={Background isolate entry point}]
@pragma('vm:entry-point')
void syncWorkerDispatcher() async {
  WidgetsFlutterBinding.ensureInitialized();
  await Firebase.initializeApp();
  // Process queue, retry pending, force sync
}
\end{lstlisting}

State machine: \texttt{idle} $\to$ \texttt{syncing} $\to$ \texttt{idle} $|$ \texttt{error}. An \texttt{\_isRunning} boolean provides coarse re-entrancy guarding.

\subsection{AIService}

Client-side gateway for AI features. Calls \texttt{aiProxy} and \texttt{ocrProcess} via \texttt{httpsCallable}. Enforces five gates in sequence:

\begin{enumerate}[leftmargin=*]
  \item \textbf{Consent:} \texttt{hasAiConsent()} / \texttt{hasOcrConsent()} from \texttt{SharedPreferences}
  \item \textbf{Cost warning:} \texttt{hasAcceptedCostWarning()}
  \item \textbf{Quota:} 50 requests/day tracked client-side (with date-rollover reset) and server-side
  \item \textbf{Connectivity:} \texttt{ConnectivityService.isOnline} must be true
  \item \textbf{Firebase availability:} \texttt{FirebaseFunctions} instance must be non-null
\end{enumerate}

If any gate fails, a localizable error string is returned (no exception thrown). The embedded system prompt enforces:
\begin{itemize}[leftmargin=*]
  \item Draft-only output (not final answers)
  \item Source grounding
  \item Academic-integrity guardrails
  \item Age-appropriate language for Pakistani Matric curriculum
\end{itemize}

\subsection{ConnectivityService}

Singleton wrapping \texttt{connectivity\_plus}. Provides:
\begin{itemize}[leftmargin=*]
  \item \texttt{isOnline} (\texttt{Future<bool>})
  \item \texttt{onConnectivityChanged} (\texttt{Stream<bool>})
\end{itemize}

Consumed by \texttt{AppShell} to render an offline banner.

\subsection{Analytics / Crashlytics / RemoteConfig}

All three follow the same pattern:
\begin{itemize}[leftmargin=*]
  \item \texttt{\_initialized} boolean guard
  \item \texttt{try/catch} on initialization
  \item Graceful degradation to no-op when Firebase is unavailable
\end{itemize}

\texttt{CrashlyticsService} is wired to \texttt{FlutterError.onError} in \texttt{main()}:
\begin{lstlisting}[language=Java]
FlutterError.onError = (details) {
  FlutterError.dumpErrorToConsole(details);
  CrashlyticsService.instance.recordError(
    details.exception,
    details.stack ?? StackTrace.current,
  );
};
\end{lstlisting}

\section{Feature Layer}

\subsection{Planner -- Greedy Knapsack Daily Budget}

The planner generates a daily study plan by:

\begin{enumerate}[leftmargin=*]
  \item Fetching all pending tasks via \texttt{getAllPendingTasks()}
  \item Scoring each task via \texttt{computeTaskScore()} (urgency $\times$ 1 + priority $\times$ 10 + fit-bonus $\times$ 5)
  \item Sorting tasks descending by score
  \item Filling the daily minute budget greedily: adding tasks in score order until $\sum \text{estimatedMinutes} \ge \text{dailyStudyMinutes}$
  \item Surfacing overflow tasks for manual rescheduling
\end{enumerate}

\textbf{Complexity:} $O(n \log n)$ due to sort, where $n$ = pending tasks. The greedy approach is not optimal for the knapsack problem but is computationally cheap and produces acceptable results for the Matric use case (small task sets).

\subsection{Focus -- Pomodoro with Android DND}

The focus screen implements a three-state machine (\texttt{idle} $\to$ \texttt{running} $\to$ \texttt{completed}). Key behaviors:

\begin{itemize}[leftmargin=*]
  \item On \texttt{dispose()} or \texttt{AppLifecycleState.paused}, an in-progress session is persisted via \texttt{\_saveSession()} to prevent data loss.
  \item Android DND is toggled via \texttt{MethodChannel('focus\_shield')}:
  \begin{itemize}[leftmargin=*]
    \item \texttt{isFocusShieldAvailable} -- checks \texttt{NotificationManager.PolicyAccess}
    \item \texttt{toggleFocusShield(enable)} -- calls \texttt{setInterruptionFilter(INTERRUPTION\_FILTER\_NONE)} or \texttt{INTERRUPTION\_FILTER\_ALL}
  \end{itemize}
  \item A foreground service (\texttt{FocusShieldForegroundService}) keeps the DND session alive against Android OS killing.
\end{itemize}

\subsection{Backup -- Encrypted Export/Import}

The backup feature follows a full DDD split:

\begin{itemize}[leftmargin=*]
  \item \textbf{Domain:} \texttt{BackupRecord}, \texttt{BackupSettings}, \texttt{RestoreResult}, \texttt{RestoreConflict} (immutable)
  \item \textbf{Domain repositories:} \texttt{BackupRepository}, \texttt{BackupEncryptionService}, \texttt{BackupStorageService}, \texttt{BackupScheduler}, \texttt{RestoreService} (abstract interfaces)
  \item \textbf{Data repositories:} \texttt{LocalBackupRepository}, \texttt{LocalBackupEncryptionService}, \texttt{LocalBackupStorageService}, \texttt{LocalRestoreService}
  \item \textbf{Application:} \texttt{BackupManager} -- composes repositories into workflows
  \item \textbf{Presentation:} \texttt{backup\_manager\_screen.dart}
\end{itemize}

\textbf{Encryption mechanics:}
\begin{itemize}[leftmargin=*]
  \item Algorithm: AES-256 via the \texttt{encrypt} package (AES-CBC under the hood)
  \item Key: 32-byte \texttt{Key.fromSecureRandom(32)}
  \item Key storage:
  \begin{itemize}[leftmargin=*]
    \item Android/iOS/macOS: \texttt{FlutterSecureStorage} (platform Keystore/Keychain)
    \item Linux/Windows: \texttt{SharedPreferences} (plaintext fallback)
  \end{itemize}
  \item File format: \texttt{MSP\_BACKUP\_V1} magic header + 16-byte IV + AES ciphertext
\end{itemize}

\textbf{Restore validation:}
\begin{enumerate}[leftmargin=*]
  \item File existence check
  \item Empty file detection
  \item Decryption failure (magic header mismatch, truncated data)
  \item Invalid JSON
  \item Missing required fields
  \item Schema version mismatch (must equal \texttt{\_schemaVersion = 2})
\end{enumerate}

Each failure mode returns a \texttt{RestoreResult} with \texttt{success: false} and an error message; no exceptions propagate to the UI.

\subsection{Revision -- SM-2 Spaced Repetition}

The revision screen displays due items where \texttt{due\_at <= now}. When a user submits confidence feedback:

\begin{enumerate}[leftmargin=*]
  \item Confidence (0--100) is mapped to SM-2 quality (0--5) via \texttt{\_confidenceToSm2Quality}
  \item \texttt{\_sm2Schedule()} computes new ease factor, interval, and repetition count
  \item The next review is scheduled at \texttt{now + interval\_days}
  \item For completed reviews with linked chapters, active task estimates are adjusted:
  \begin{align}
  \text{ratio} &= \frac{\text{originalEstimate} > 0}{100.0} : 0.5 \\
  \text{newEstimate} &= \text{clamp}(\text{round}(\text{originalEstimate} \cdot \text{ratio}),\ 1,\ \text{originalEstimate} \cdot 3)
  \end{align}
\end{enumerate}

\section{Cloud Layer}

\subsection{Cloud Functions}

Four Firebase Cloud Functions (TypeScript, Firebase Functions v2):

\begin{table}[H]
\centering
\caption{Firebase Cloud Functions}
\label{tab:cloud-functions}
\begin{tabular}{llp{7cm}}
\toprule
\textbf{Function} & \textbf{Trigger} & \textbf{Purpose} \\
\midrule
\texttt{aiProxy} & \texttt{https.onCall} & Proxies AI requests to OpenAI/Anthropic/Gemini with consent + quota checks \\
\texttt{ocrProcess} & \texttt{https.onCall} & Enqueues OCR jobs via Google Vision API with duplicate suppression \\
\texttt{syncWorker} & \texttt{pubsub.schedule('every 15 minutes')} & Processes \texttt{sync\_outbox}, applying mutations to target collections \\
\texttt{pruneTombstones} & \texttt{pubsub.schedule('every 24 hours')} & Deletes \texttt{sync\_tombstones} older than 30 days \\
\bottomrule
\end{tabular}
\end{table}

\paragraph{\texttt{aiProxy} flow:}
\begin{enumerate}[leftmargin=*]
  \item Validate \texttt{context.auth.uid} exists
  \item Read \texttt{ai\_consents/\{userId\}} document
  \item Enforce 50-request/day quota by querying \texttt{ai\_requests} for current day
  \item Record request in \texttt{ai\_requests} (prompt truncated to 2000 chars)
  \item Resolve provider/model, call appropriate AI SDK
  \item Validate response non-empty and $\le$ 4000 chars
  \item Return \texttt{\{requestId, provider, model, output, notice\}}
\end{enumerate}

\paragraph{\texttt{ocrProcess} flow:}
\begin{enumerate}[leftmargin=*]
  \item Validate auth and OCR consent
  \item Duplicate detection: query \texttt{ocr\_jobs} for matching \texttt{userId} + \texttt{filePath} with status \texttt{queued}/\texttt{processing} within last hour
  \item Create \texttt{ocr\_jobs} document
  \item Invoke \texttt{ImageAnnotatorClient.asyncBatchAnnotateFiles()} for PDF text detection
  \item Return job ID for client polling
\end{enumerate}

\paragraph{\texttt{syncWorker} flow:}
\begin{enumerate}[leftmargin=*]
  \item Read \texttt{sync\_outbox} where \texttt{synced == false} (limit 100)
  \item For each document, resolve target collection reference
  \item Apply \texttt{delete}, \texttt{update}, or \texttt{set(..., \{merge: true\})}
  \item Mark document \texttt{synced: true} with \texttt{syncedAt} timestamp
  \item On failure, write error to document's \texttt{error} field; function does NOT throw
\end{enumerate}

\paragraph{\texttt{pruneTombstones} flow:}
\begin{enumerate}[leftmargin=*]
  \item Query \texttt{sync\_tombstones} where \texttt{deletedAt < cutoff} (30 days ago, limit 500)
  \item Batch delete
  \item Commit
\end{enumerate}

\subsection{Firestore Security Rules}

The rules implement a strict \textbf{owner-based access control} model:

\begin{verbatim}
function isOwner(userId) {
  return isAuthenticated() && request.auth.uid == userId;
}
\end{verbatim}

\begin{table}[H]
\centering
\caption{Firestore collection access matrix}
\label{tab:firestore-rules}
\begin{tabular}{lll}
\toprule
\textbf{Collection} & \textbf{Read} & \textbf{Write/Create} \\
\midrule
\texttt{sync\_outbox} & Authenticated + owner & Authenticated + owner \\
\texttt{sync\_tombstones} & Authenticated + owner & Authenticated + owner \\
\texttt{sync\_conflicts} & Authenticated + owner & Authenticated + owner \\
\texttt{user\_profiles} & Owner & Owner \\
\texttt{quizzes} & Authenticated & Owner \\
\texttt{quiz\_attempts} & Owner & Owner \\
\bottomrule
\end{tabular}
\end{table}

Collections \texttt{ai\_consents}, \texttt{ai\_requests}, \texttt{ocr\_jobs}, and \texttt{ai\_hallucination\_reports} are not explicitly listed because they are accessed exclusively by Cloud Functions (which bypass client-side rules) or by client methods with embedded ownership checks.

\subsection{Storage Rules}

\begin{table}[H]
\centering
\caption{Firebase Storage access rules}
\label{tab:storage-rules}
\begin{tabular}{ll}
\toprule
\textbf{Path} & \textbf{Access} \\
\midrule
\texttt{users/\{userId\}/**}& Read/Write: owner only \\
\texttt{ai-uploads/\{userId\}/**}& Read/Write: owner only \\
\texttt{ocr-output/\{userId\}/**}& Read/Write: owner only \\
\texttt{study-groups/\{groupId\}/**}& Read: authenticated; Write: denied \\
\bottomrule
\end{tabular}
\end{table}

% ═══════════════════════════════════════════════════════════════════════════════
\chapter{Memory Management \& Hardware/Resource Interaction}
% ═══════════════════════════════════════════════════════════════════════════════

\section{Dart Isolate Model}

The Flutter application operates primarily in a \textbf{single Dart isolate} with a single-threaded event loop. Database operations are \texttt{async} and yield to the event loop at \texttt{await} points; they are not dispatched to separate OS threads.

\begin{itemize}[leftmargin=*]
  \item \textbf{sqflite path:} Queries execute on the main isolate's event loop. SQLite itself is thread-safe at the C level, but the Dart wrapper serializes access through the singleton \texttt{DatabaseHelper}.
  \item \textbf{Drift path:} \texttt{NativeDatabase.createInBackground()} offloads SQLite I/O to a \textbf{background Dart isolate} using \texttt{Isolate.spawn}. This prevents I/O latency from blocking the UI thread.
  \item \textbf{workmanager:} Spawns a \textbf{separate Dart isolate} for background sync. The entry point is annotated with \texttt{@pragma('vm:entry-point')} to prevent tree-shaking in release builds. This isolate initializes its own \texttt{Firebase.initializeApp()} instance.
\end{itemize}

\section{Memory Layout \& Allocation Profiles}

\begin{table}[H]
\centering
\caption{Memory allocation profiles}
\label{tab:memory}
\begin{tabular}{llll}
\toprule
\textbf{Resource} & \textbf{Location} & \textbf{Lifetime} & \textbf{Notes} \\
\midrule
SQLite database file & \texttt{getApplicationDocumentsDirectory()/study\_planner.db} & Persistent & ~11 tables; estimated 2--10 MB for typical Matric student data \\
Encryption key (backup) & \texttt{FlutterSecureStorage} or \texttt{SharedPreferences} & Persistent & 32 bytes; plaintext fallback on desktop \\
Backup files & \texttt{documents/backups/} & Persistent & AES-256 ciphertext; user-initiated \\
Sync outbox (Firestore) & Cloud & Ephemeral & Marked \texttt{synced: true} after processing \\
Offline queue (SharedPreferences) & Local & Ephemeral & Max 3 retry attempts per operation \\
\texttt{sync\_meta} table & SQLite & Persistent & Local outbox ledger; entries accumulate until synced \\
Focus session state & In-memory + SQLite & Session + persistent & Saved on \texttt{dispose()} to prevent data loss \\
\bottomrule
\end{tabular}
\end{table}

\section{Zero-Copy Abstractions}

The application does not employ explicit zero-copy abstractions. Data flows through:
\begin{enumerate}[leftmargin=*]
  \item SQLite returns \texttt{Map<String, dynamic>} rows
  \item \texttt{fromMap()} constructs model instances (copy)
  \item \texttt{toMap()} serializes back to \texttt{Map<String, dynamic>}
  \item Firestore writes receive these maps directly
\end{enumerate}

For large payloads (AI prompts, OCR PDFs), data is transferred as:
\begin{itemize}[leftmargin=*]
  \item AI prompts: $\le$2000 chars truncated server-side
  \item AI responses: $\le$4000 chars validated server-side
  \item OCR files: Uploaded to Firebase Storage; processing is asynchronous via polling
\end{itemize}

\section{Native Android Integration}

\paragraph{FocusShieldForegroundService (Kotlin):}
\begin{lstlisting}[language=Java,caption={FocusShieldForegroundService (Kotlin)}]
class FocusShieldForegroundService : Service() {
  override fun onStartCommand(intent: Intent, flags: Int, startId: Int): Int {
    val notification = createNotification("Focus session active", "DND mode enabled")
    startForeground(NOTIFICATION_ID, notification)
    // Toggle DND via NotificationManager
    return START_STICKY
  }
}
\end{lstlisting}

The foreground service holds a persistent notification and prevents the OS from killing the DND toggle process during active focus sessions.

\paragraph{FocusAccessibilityService (Kotlin):} An optional accessibility service for usage-time tracking (not actively used in the current feature set).

\paragraph{App Shortcuts:} \texttt{MainActivity.kt} registers the \texttt{focus\_shield} platform channel and handles the \texttt{app\_shortcuts} channel for launch-intent detection.

\section{Platform-Specific Constraints}

\begin{table}[H]
\centering
\caption{Platform capability matrix}
\label{tab:platforms}
\begin{tabular}{lllll}
\toprule
\textbf{Platform} & \textbf{sqflite} & \textbf{Drift} & \textbf{Secure Storage} & \textbf{Notes} \\
\midrule
Android & \checkmark & \checkmark (default) & Keystore & Primary target \\
iOS & \checkmark & \texttimes & Keychain & Not primary target \\
Linux & \checkmark (FFI) & \texttimes & SharedPreferences (plaintext) & Desktop support \\
Windows & \checkmark (FFI) & \texttimes & SharedPreferences (plaintext) & Desktop support \\
macOS & \checkmark (FFI) & \texttimes & Keychain & Desktop support \\
Web & \texttimes & \texttimes & \texttimes & Explicitly unsupported \\
\bottomrule
\end{tabular}
\end{table}

% ═══════════════════════════════════════════════════════════════════════════════
\chapter{Data Flow, Interfaces \& API Specifications}
% ═══════════════════════════════════════════════════════════════════════════════

\section{End-to-End Data Flow: Local Write $\to$ Cloud Sync}

\paragraph{Pathway A: Local Mutation $\to$ Cloud Synchronization}

\begin{enumerate}[leftmargin=*]
  \item \textbf{UI Invocation:} User action (e.g., complete chapter) $\to$ UI calls \texttt{DatabaseRepository.updateChapterStatus(id, 'completed')}
  \item \textbf{Local Database Write:} \texttt{DatabaseHelper.updateChapterStatus()} $\to$ SQLite: \texttt{UPDATE chapters SET status='completed', completed\_at=... WHERE id=?}
  \item \textbf{Sync Metadata Tracking (INV-2):} \texttt{DatabaseRepository.\_trackChange('chapters', id.toString(), 'update', \{...\})} $\to$ SQLite: \texttt{INSERT OR REPLACE INTO sync\_meta ...}
  \item \textbf{Online Cloud Enqueue:} If \texttt{ConnectivityService.isOnline == true}: \texttt{SyncService.enqueueLocalChange(userId, 'chapters', '42', 'update', \{...\})} $\to$ Firestore: \texttt{sync\_outbox.add(\{userId, entity, entityId, operation, data, synced: false\})}
  \item \textbf{Background Cloud Processing (every 15 min):} Cloud Function \texttt{syncWorker} queries \texttt{sync\_outbox} where \texttt{synced == false}, applies mutations, marks \texttt{synced: true}
  \item \textbf{Client-Side Fallback:} \texttt{SyncService.retryPending()} reads \texttt{sync\_outbox}, applies create/update/delete to target collection, marks \texttt{synced: true}
  \item \textbf{Offline Queue (if Step 4 fails):} \texttt{OfflineQueueService.enqueue(\{type: 'sync', payload: \{...\}\})} $\to$ \texttt{SharedPreferences}; Workmanager (15 min, \texttt{NetworkType.connected}) processes queue $\to$ \texttt{SyncService.enqueueLocalChange()}
\end{enumerate}

\section{End-to-End Data Flow: AI Request}

\paragraph{Pathway B: AI-Assisted Task Breakdown}

\begin{enumerate}[leftmargin=*]
  \item \textbf{UI Invocation:} User taps ``AI Task Breakdown'' on a chapter $\to$ UI calls \texttt{AIService.generateTaskBreakdown(chapterTitle)}
  \item \textbf{Gate Checks (sequential, short-circuit on failure):}
  \begin{itemize}[leftmargin=*]
    \item \texttt{hasAiConsent()} $\to$ \texttt{SharedPreferences}
    \item \texttt{hasAcceptedCostWarning()} $\to$ \texttt{SharedPreferences}
    \item \texttt{getQuotaRemaining()} $\to$ \texttt{SharedPreferences} (date-rollover)
    \item \texttt{ConnectivityService.isOnline} $\to$ \texttt{Future<bool>}
    \item Any gate fails $\to$ return localized error string
  \end{itemize}
  \item \textbf{Client-Side Quota Increment:} \texttt{SharedPreferences}: increment \texttt{ai\_requests\_today}, set \texttt{ai\_requests\_date}
  \item \textbf{Cloud Function Call:} \texttt{FirebaseFunctions.httpsCallable('aiProxy')}
  \begin{itemize}[leftmargin=*]
    \item Server-side: Auth check, consent check, quota check, record request, resolve provider/model, call AI SDK, validate output
    \item Returns \texttt{\{requestId, provider, model, output, notice\}}
  \end{itemize}
  \item \textbf{UI Display:} UI renders output with ``AI responses may contain errors'' disclaimer
\end{enumerate}

\section{End-to-End Data Flow: Backup Creation}

\paragraph{Pathway C: Encrypted Local Backup}

\begin{enumerate}[leftmargin=*]
  \item \textbf{BackupManager.createBackup():} \texttt{DatabaseRepository} collects all entities; serializes to JSON
  \item \textbf{Encryption:} \texttt{LocalBackupEncryptionService.encryptData(jsonString)}
  \begin{itemize}[leftmargin=*]
    \item Retrieves/generates 32-byte AES key from \texttt{FlutterSecureStorage}
    \item Generates 16-byte random IV
    \item AES-CBC encrypt: \texttt{ciphertext = AES.encrypt(pkcs7\_pad(json), key, iv)}
    \item Returns: \texttt{"MSP\_BACKUP\_V1"} + \texttt{iv\_base64} + \texttt{ciphertext\_base64}
  \end{itemize}
  \item \textbf{File Write:} \texttt{LocalBackupStorageService.writeBackupFile(encryptedData)} $\to$ Path: \texttt{getApplicationDocumentsDirectory()/backups/backup\_<timestamp>.msp}
  \item \textbf{Metadata Record:} \texttt{DatabaseHelper.insertBackupRecord(destination, status)} $\to$ SQLite: \texttt{INSERT INTO backup\_records ...}
\end{enumerate}

\section{Public API Specifications}

\subsection{DatabaseRepository Interface}

\begin{lstlisting}[language=Java,caption={DatabaseRepository public interface}]
class DatabaseRepository {
  // User
  Future<UserModel?> getUserProfile();
  Future<int> createUserProfile(UserModel profile);
  Future<int> updateUserProfile(UserModel profile);

  // Subjects
  Future<List<SubjectModel>> getSubjects();
  Future<int> insertSubject(SubjectModel subject);
  Future<int> deleteSubject(int id);

  // Chapters
  Future<List<ChapterModel>> getChaptersBySubject(int subjectId);
  Future<int> insertChapter(ChapterModel chapter);
  Future<int> updateChapterStatus(int id, String status);
  Future<void> setChapterCompletion(int chapterId, bool completed);
  Future<ChapterModel?> getChapterById(int id);

  // Tasks
  Future<List<TaskModel>> getTodayTasks();
  Future<List<TaskModel>> getAllPendingTasks();
  Future<List<TaskModel>> getTasksForDate(DateTime date);
  Future<List<TaskModel>> getTasksForDateRange(DateTime start, DateTime end);
  Future<int> insertTask(TaskModel task);
  Future<int> updateTask(int id, Map<String, dynamic> data);
  Future<int> deleteTask(int id);
  int computeTaskScore(TaskModel task, int dailyMinutes);

  // Focus Sessions
  Future<FocusSessionModel?> getActiveFocusSession();
  Future<int> insertFocusSession(FocusSessionModel session);
  Future<int> updateFocusSession(int id, Map<String, dynamic> data);

  // Revision (SM-2)
  Future<List<RevisionItemModel>> getDueRevisions();
  Future<void> recordRevisionFeedback(int revisionId, int confidence, String status);

  // Analytics
  Future<List<Map<String, dynamic>>> getRecallTrends();
  Future<List<Map<String, dynamic>>> getSubjectConfidenceDecay();
  Future<List<Map<String, dynamic>>> getTaskEstimateAccuracy();
  Future<List<Map<String, dynamic>>> getProductiveTimeInsights();
  Future<List<Map<String, dynamic>>> getMissedDayPatterns();
  Future<List<Map<String, dynamic>>> getSubjectEstimateAccuracy();

  // Search
  Future<List<dynamic>> search(String query);

  // Backup
  Future<List<Map<String, dynamic>>> getBackupHistory();
  String exportAnkiCsv(List<SubjectModel> subjects, List<ChapterModel> chapters);
  List<Map<String, dynamic>> importAnkiCsv(String csvContent, int subjectId);
}
\end{lstlisting}

\subsection{SyncService Interface}

\begin{lstlisting}[language=Java,caption={SyncService public interface}]
class SyncService {
  Future<void> enqueueLocalChange({
    required String userId,
    required String entity,
    required String entityId,
    required String operation, // 'create' | 'update' | 'delete'
    required Map<String, dynamic> data,
  });

  Future<void> markEntityDeleted({
    required String userId,
    required String entity,
    required String entityId,
  });

  Future<List<Map<String, dynamic>>> getPendingOutbox(String userId);
  Future<void> markSynced(String outboxId);
  Future<void> recordConflict({
    required String userId,
    required String entity,
    required String entityId,
    required Map<String, dynamic> localData,
    required Map<String, dynamic> remoteData,
  });
  Future<void> retryPending();
  Future<int> pruneTombstones();
  Future<int> getPendingOutboxCount();
  Future<int> getPendingConflictsCount();
}
\end{lstlisting}

\subsection{AIService Interface}

\begin{lstlisting}[language=Java,caption={AIService public interface}]
class AIService {
  Future<String> generateTaskBreakdown(String chapterTitle, String subjectName);
  Future<String> generateRevisionDraft(String chapterTitle);
  Future<String> generateFlashcards(String chapterTitle);
  Future<String> generateQuiz(String chapterTitle, int questionCount);
  Future<String> answerSubjectQuestion(String question, String context);
  Future<String> processOcrScan(String filePath, String language);
}
\end{lstlisting}

All methods return \texttt{String} -- either the AI output or a localizable error message. No exceptions are thrown for gating failures.

\subsection{Type System \& Invariants}

The Dart type system is used extensively but relies on runtime checks for external data (SQLite rows, Firestore documents, JSON). Key type invariants:

\begin{itemize}[leftmargin=*]
  \item \texttt{TaskModel.status} is one of: \texttt{'pending'}, \texttt{'in\_progress'}, \texttt{'completed'}
  \item \texttt{FocusSessionModel.status} is one of: \texttt{'running'}, \texttt{'completed'}, \texttt{'paused'}
  \item \texttt{revision\_items.recall\_confidence} is \texttt{int} in range [0, 100]
  \item \texttt{revision\_items.ease\_factor} is \texttt{double} $\ge$ 1.3
  \item \texttt{sync\_meta.sync\_status} is one of: \texttt{'pending'}, \texttt{'synced'}, \texttt{'conflict'}
  \item \texttt{chapters.content\_tier} is one of: \texttt{'official'}, \texttt{'supplementary'}
\end{itemize}

The \texttt{sync\_meta} table has a \texttt{UNIQUE(entity, local\_id)} constraint enforced at the SQLite level, ensuring that \texttt{\_trackChange()} calls are idempotent for the same entity mutation.

% ═══════════════════════════════════════════════════════════════════════════════
\chapter{Verification, Safety \& Security Model}
% ═══════════════════════════════════════════════════════════════════════════════

\section{Static Analysis \& Testing}

\begin{table}[H]
\centering
\caption{Verification toolchain}
\label{tab:verification}
\begin{tabular}{lll}
\toprule
\textbf{Layer} & \textbf{Tool} & \textbf{Coverage} \\
\midrule
Dart/Flutter & \texttt{flutter analyze} (\texttt{flutter\_lints} baseline) & All Dart code \\
Unit tests & \texttt{flutter test} & 21 unit test files + widget tests \\
Firestore rules & Jest + \texttt{@firebase/rules-unit-testing} & Rule logic validation \\
Cloud Functions & Jest & Function logic + error paths \\
CI/CD & GitHub Actions & Format, analyze, test, build on push/PR \\
\bottomrule
\end{tabular}
\end{table}

Test files cover: AI service, sync service, offline queue, backup encryption/restore, planner logic, repository, focus screen, settings, legal screen, OCR service, production fixes, and sync worker.

\section{Defensive Programming Guarantees}

\paragraph{Type Safety in Data Models.}
Every \texttt{fromMap()} constructor uses \texttt{is} checks before casting:
\begin{lstlisting}[language=Java]
studentName: map['student_name'] is String ? map['student_name'] as String : '',
\end{lstlisting}
This protects against partial rows, legacy schemas, and malformed backup imports.

\paragraph{Firebase Error Swallowing.}
All Firebase-dependent services catch \texttt{FirebaseException} and return safe defaults. This is an intentional invariant -- no Firebase dependency may crash the app. Verified by tests in \texttt{test/unit/sync\_service\_test.dart}.

\paragraph{\texttt{FlutterError.onError}.}
Captures all Flutter framework errors and routes them to Crashlytics when available:
\begin{lstlisting}[language=Java]
FlutterError.onError = (details) {
  FlutterError.dumpErrorToConsole(details);
  CrashlyticsService.instance.recordError(details.exception, details.stack);
};
\end{lstlisting}

\section{Security Architecture}

\subsection{Threat Model}

\begin{table}[H]
\centering
\caption{Security threat model and mitigations}
\label{tab:threat-model}
\small
\begin{tabular}{>{\bfseries}l p{8cm}}
\toprule
\textbf{Threat} & \textbf{Mitigation} \\
\midrule
Unauthorized cloud data access & Firestore/Storage owner-based rules; all collections require \texttt{request.auth.uid == userId} \\
Data resurrection after delete & 30-day tombstone retention in \texttt{sync\_tombstones}; Cloud Function applies tombstones before allowing new creates \\
AI quota abuse & Client-side (50/day SharedPreferences) + server-side (Firestore \texttt{ai\_requests} count) enforcement \\
AI consent bypass & Both client and server enforce \texttt{ai\_consents/\{userId\}.aiAssistance} \\
Backup decryption by attacker & AES-256 encryption; key in platform Keystore/Keychain (not plaintext on mobile) \\
Backup tampering & Magic header validation (\texttt{MSP\_BACKUP\_V1}); schema version check on restore \\
Unauthorized DND toggle & Android \texttt{ACCESS\_NOTIFICATION\_POLICY} permission required; app directs user to system settings if missing \\
Sync conflict data loss & Last-write-wins for non-deletes; true conflicts recorded in \texttt{sync\_conflicts} for manual resolution \\
Excessive AI cost & 50-request/day quota; cost-warning acceptance required before first AI use \\
OCR duplicate processing & Server-side duplicate detection: same \texttt{userId} + \texttt{filePath} + status within 1 hour \\
\bottomrule
\end{tabular}
\end{table}

\subsection{Memory Safety Invariants}

\begin{itemize}[leftmargin=*]
  \item All SQLite operations use parameterized queries (no string concatenation in \texttt{where} clauses)
  \item Foreign keys use \texttt{ON DELETE CASCADE} / \texttt{ON DELETE SET NULL} to prevent orphaned rows
  \item The \texttt{sync\_meta} \texttt{UNIQUE(entity, local\_id)} constraint prevents duplicate outbox entries
  \item AES encryption uses random IV per operation (no ECB mode; AES-CBC via \texttt{encrypt} package)
  \item Encryption keys are never logged or transmitted
\end{itemize}

\subsection{Permission \& Consent Model}

\begin{table}[H]
\centering
\caption{Consent and permission model}
\label{tab:consent}
\begin{tabular}{lll}
\toprule
\textbf{Feature} & \textbf{Client Gate} & \textbf{Server Gate} \\
\midrule
Cloud sync & None (always enabled when online) & Firestore rules (authenticated) \\
AI assistance & \texttt{hasAiConsent()} + \texttt{hasAcceptedCostWarning()} + quota & \texttt{ai\_consents/\{userId\}.aiAssistance} \\
OCR scanning & \texttt{hasOcrConsent()} + connectivity & \texttt{ai\_consents/\{userId\}.ocrScanning} \\
Analytics & \texttt{AnalyticsService.initialize()} (opt-out via settings) & N/A (Firebase config) \\
Crashlytics & \texttt{CrashlyticsService.initialize()} & N/A (Firebase config) \\
Android DND & \texttt{ACCESS\_NOTIFICATION\_POLICY} runtime permission & N/A (OS-level) \\
\bottomrule
\end{tabular}
\end{table}

\subsection{Licensing \& Governance}

\begin{itemize}[leftmargin=*]
  \item \textbf{License:} Apache License, Version 2.0
  \item \textbf{Copyright:} CodeSym, 2026
  \item \textbf{Third-party dependencies:} Managed via \texttt{pubspec.yaml} and \texttt{package.json}; all licenses are permissive (MIT, BSD, Apache) compatible with the project's Apache 2.0 license
  \item \textbf{AI provider terms:} Users are responsible for compliance with OpenAI/Anthropic/Gemini/OpenRouter terms of service; the app does not store API keys client-side (keys are in Firebase Functions config/environment variables)
  \item \textbf{Data retention:} Documented in \texttt{docs/legal/data\_retention\_policy.md}; sync tombstones pruned after 30 days
\end{itemize}

% ═══════════════════════════════════════════════════════════════════════════════
\chapter{Performance Profiles \& Benchmarking}
% ═══════════════════════════════════════════════════════════════════════════════

\section{Theoretical Performance Bounds}

\begin{table}[H]
\centering
\caption{Algorithmic complexity of primary operations}
\label{tab:complexity}
\begin{tabular}{lll}
\toprule
\textbf{Operation} & \textbf{Complexity} & \textbf{Notes} \\
\midrule
\texttt{getTodayTasks()} & $O(n)$ & Full table scan with date filter; $n$ = total tasks \\
\texttt{getAllPendingTasks()} & $O(n)$ & Full table scan with status filter \\
\texttt{getTasksForDate(date)} & $O(\log n)$ & Assumes index on \texttt{scheduled\_at} \\
\texttt{getSubjectProgress()} & $O(s + c)$ & Join over subjects + chapters; $s$ = subjects, $c$ = chapters \\
\texttt{getRecallTrends()} & $O(r)$ & Raw query with GROUP BY; $r$ = revision items \\
\texttt{computeTaskScore(task, dailyMinutes)} & $O(1)$ & Pure arithmetic \\
\texttt{\_sm2Schedule(...)} & $O(1)$ & Pure arithmetic \\
Greedy knapsack (planner) & $O(n \log n)$ & Sort + linear scan \\
\texttt{retryPending()} & $O(o)$ & $o$ = outbox items (capped at 100) \\
\texttt{search(query)} & $O(n \times m)$ & $n$ = total rows across 5 tables; $m$ = avg row size \\
\bottomrule
\end{tabular}
\end{table}

\section{SQLite Query Optimization}

The application uses raw SQL for analytics queries with explicit JOINs and GROUP BY. The schema lacks explicit indexes beyond primary keys. For typical Matric data volumes ($\le$100 subjects, $\le$1000 chapters, $\le$5000 tasks), query times are sub-10ms on modern mobile hardware.

Recommended index additions for scaling:
\begin{verbatim}
CREATE INDEX idx_study_tasks_scheduled_at ON study_tasks(scheduled_at);
CREATE INDEX idx_study_tasks_status ON study_tasks(status);
CREATE INDEX idx_revision_items_due_at ON revision_items(due_at);
CREATE INDEX idx_focus_sessions_task_id ON focus_sessions(task_id);
CREATE INDEX idx_sync_meta_entity_local_id ON sync_meta(entity, local_id);
\end{verbatim}

\section{Background Sync Efficiency}

\begin{itemize}[leftmargin=*]
  \item \textbf{Frequency:} 15 minutes (configurable via Remote Config)
  \item \textbf{Batch size:} 100 outbox items per \texttt{syncWorker} invocation
  \item \textbf{Tombstone pruning:} 500 tombstones per \texttt{pruneTombstones} invocation (daily)
  \item \textbf{Network constraint:} \texttt{NetworkType.connected} -- sync does not run on metered/unmetered ambiguity; requires active connectivity
\end{itemize}

The 15-minute interval is a trade-off between cloud consistency and battery/network usage. For a student population primarily on mobile data, this minimizes background network activity while maintaining reasonable sync freshness.

\section{Offline Queue Bounds}

\begin{itemize}[leftmargin=*]
  \item \textbf{Max retries per operation:} 3
  \item \textbf{Persistence:} \texttt{SharedPreferences} (JSON array)
  \item \textbf{Eviction:} Operations exceeding max retries are silently dropped
  \item \textbf{Memory footprint:} Negligible; queue is loaded into memory on access and written back on modification
\end{itemize}

\section{Focus Session Timer Precision}

The Pomodoro timer uses Dart's \texttt{Timer.periodic()} with a 1-second tick. On \texttt{AppLifecycleState.paused}, the session is persisted to SQLite. On resume, the timer continues from its last tick. There is no compensation for time spent in background; the timer counts wall-clock seconds.

\section{AI Request Throughput}

\begin{itemize}[leftmargin=*]
  \item \textbf{Client-side quota:} 50 requests/day per user
  \item \textbf{Server-side quota:} Same 50-request/day limit enforced in \texttt{aiProxy}
  \item \textbf{Prompt truncation:} 2000 characters (server-side)
  \item \textbf{Response cap:} 4000 characters (validated server-side)
  \item \textbf{Timeout:} Inherited from Firebase Functions default (60 seconds for \texttt{https.onCall})
\end{itemize}

\section{Cache Efficiency}

The application does not implement an explicit caching layer beyond SQLite. The \texttt{DatabaseHelper} singleton caches the open \texttt{Database} connection, avoiding repeated file open/close overhead. Riverpod providers cache derived state (e.g., \texttt{todayTasksProvider} caches the result until dependencies change).

% ═══════════════════════════════════════════════════════════════════════════════
\chapter{Deployment, Integration \& Future Expansion Roadmap}
% ═══════════════════════════════════════════════════════════════════════════════

\section{Production Deployment Requirements}

\begin{table}[H]
\centering
\caption{Deployment targets and artifacts}
\label{tab:deployment}
\begin{tabular}{llll}
\toprule
\textbf{Component} & \textbf{Target} & \textbf{Build Command} & \textbf{Artifact} \\
\midrule
Android App & Android 5.0+ (API 21+) & \texttt{flutter build apk --release} or \texttt{flutter build appbundle} & \texttt{.apk} or \texttt{.aab} \\
Linux Desktop & x64 Linux & \texttt{flutter build linux} & Binary in \texttt{build/linux/x64/release/bundle/} \\
Firebase Functions & Node.js 20+ & \texttt{cd functions \&\& npm run deploy} & Cloud Functions \\
Firestore Rules & -- & \texttt{firebase deploy --only firestore:rules} & Security rules \\
Storage Rules & -- & \texttt{firebase deploy --only storage} & Security rules \\
\bottomrule
\end{tabular}
\end{table}

\section{Build Toolchain}

\begin{verbatim}
+---------------------------------------------------+
|                 BUILD PIPELINE                    |
|                                                   |
|  flutter analyze                                   |
|      |                                            |
|      v                                            |
|  flutter test                                      |
|      |                                            |
|      v                                            |
|  flutter build apk --release (Android)             |
|  flutter build linux (Desktop)                     |
|      |                                            |
|      v                                            |
|  cd functions && npm ci && npm run build          |
|      |                                            |
|      v                                            |
|  firebase deploy --only functions,firestore,storage|
+---------------------------------------------------+
\end{verbatim}

\section{CI/CD (GitHub Actions)}

\texttt{.github/workflows/ci.yml} triggers on push/PR to \texttt{main}/\texttt{develop}:
\begin{enumerate}[leftmargin=*]
  \item \texttt{flutter pub get}
  \item \texttt{dart format --set-exit-if-changed .}
  \item \texttt{flutter analyze}
  \item \texttt{flutter test}
  \item Build Android AAB (release) or APK (debug on PR)
  \item Build Linux release binary
\end{enumerate}

\texttt{.github/workflows/firestore-rules.yml} triggers on changes to rules:
\begin{enumerate}[leftmargin=*]
  \item \texttt{cd functions \&\& npm ci}
  \item \texttt{npm run test:rules}
\end{enumerate}

\section{Extensibility Mechanisms}

\subsection{Feature Flags}

\texttt{lib/core/feature\_flags.dart} defines compile-time and runtime feature flags:

\begin{lstlisting}[language=Java,caption={Feature flags definition}]
class FeatureFlags {
  static const bool enableAI = bool.fromEnvironment('ENABLE_AI', defaultValue: true);
  static const bool enableOCR = bool.fromEnvironment('ENABLE_OCR', defaultValue: true);
  static const bool enableSpacedRepetition = bool.fromEnvironment('ENABLE_SR', defaultValue: true);
  static const bool enableEncryptedStorage = bool.fromEnvironment('ENABLE_ENCRYPTED_STORAGE', defaultValue: true);
  // ...
}
\end{lstlisting}

Flags can be overridden at build time via \texttt{--dart-define=ENABLE\_AI=false}.

\subsection{DDD Pattern for Complex Features}

The backup feature demonstrates the full domain/data/presentation/application split. New complex features should follow this pattern:
\begin{itemize}[leftmargin=*]
  \item Define abstract interfaces in \texttt{domain/repositories/}
  \item Implement in \texttt{data/repositories/}
  \item Compose workflows in \texttt{application/}
  \item Render in \texttt{presentation/}
\end{itemize}

This enables mock-based testing without UI and allows alternative implementations (e.g., cloud backup) by implementing the same interfaces.

\subsection{AI Provider Extensibility}

The \texttt{aiProxy} Cloud Function supports multiple AI providers via environment variables:
\begin{itemize}[leftmargin=*]
  \item \texttt{AI\_PROVIDER} -- selects provider (\texttt{openai}, \texttt{anthropic}, \texttt{gemini}, \texttt{openrouter})
  \item \texttt{AI\_API\_KEY} -- OpenAI/OpenRouter API key
  \item \texttt{ANTHROPIC\_API\_KEY} -- Anthropic API key
  \item \texttt{GEMINI\_API\_KEY} -- Google Gemini API key
\end{itemize}

Adding a new provider requires:
\begin{enumerate}[leftmargin=*]
  \item Adding a client factory in \texttt{functions/src/index.ts}
  \item Adding a branch in the \texttt{aiProxy} try block
  \item Adding the provider to \texttt{providerDefaults}
\end{enumerate}

\subsection{Syllabus Template System}

The \texttt{syllabus\_templates} table allows importing/exporting structured syllabus data as JSON. The system seeds default Pakistani Matric subjects on first launch if no templates exist. New boards or class levels can be supported by:
\begin{enumerate}[leftmargin=*]
  \item Adding template entries to \texttt{MatricSubjects.subjects}
  \item Or importing JSON templates with \texttt{\{subjects: [...], chapters: [...]\}} structure
\end{enumerate}

\section{Architectural Roadmap}

\begin{table}[H]
\centering
\caption{Phased architectural roadmap}
\label{tab:roadmap}
\begin{tabular}{llp{8cm}}
\toprule
\textbf{Phase} & \textbf{Milestone} & \textbf{Technical Work} \\
\midrule
v1.1 & Multi-device sync hardening & Implement vector clocks for causality tracking; replace last-write-wins with CRDT merge for task titles \\
v1.2 & Conflict resolution UX & Build guided merge UI for \texttt{sync\_conflicts}; auto-archive resolved conflicts \\
v1.3 & Performance scaling & Add composite indexes on \texttt{scheduled\_at}, \texttt{status}, \texttt{due\_at}; implement SQLite FTS5 for full-text search \\
v2.0 & Cloud-native backup & Implement \texttt{CloudBackupRepository} interface; store encrypted backups in Firebase Storage with user-controlled keys \\
v2.1 & Collaborative study groups & Extend Firestore rules for group-scoped collections; implement real-time presence via Firestore listeners \\
v2.2 & Adaptive ML planning & Replace greedy knapsack with reinforcement learning policy; train on user completion patterns \\
v3.0 & Platform expansion & Add iOS/macOS support via \texttt{flutter\_secure\_storage} Keychain integration; enable Drift on all platforms \\
\bottomrule
\end{tabular}
\end{table}

\section{Known Technical Debt}

\begin{table}[H]
\centering
\caption{Technical debt register}
\label{tab:debt}
\begin{tabular}{lll}
\toprule
\textbf{Item} & \textbf{Severity} & \textbf{Mitigation} \\
\midrule
Drift schema maintained in parallel with raw SQLite & Medium & Consider migrating fully to Drift and dropping \texttt{DatabaseHelper} \\
\texttt{sync\_meta} table not exposed via Drift DAOs & Low & Add \texttt{SyncMeta} Drift table + DAO \\
Plaintext encryption key fallback on desktop & Medium & Integrate with OS keyrings (libsecret on Linux, Credential Manager on Windows) \\
No FTS5 for search & Low & Add FTS5 virtual tables for \texttt{study\_tasks}, \texttt{chapters}, \texttt{subjects} \\
\texttt{retryPending()} last-write-wins & Medium & Implement server-side version vectors or operational transforms \\
Hard-coded 15-minute sync interval & Low & Expose via Remote Config + user settings \\
\bottomrule
\end{tabular}
\end{table}

% ═══════════════════════════════════════════════════════════════════════════════
\chapter*{Appendix A: ASCII Data Flow Summary}
% ═══════════════════════════════════════════════════════════════════════════════

\begin{blockverbatim}
+-------------------------------------------------------------+
|                     COMPLETE SYSTEM VIEW                    |
+-------------------------------------------------------------+
|                                                             |
|  +----------+    Riverpod     +--------------+ sqflite/Drift|
|  |  Flutter  | ------------>  | DatabaseRepo | ------------>|
|  |   UI     | <------------  |  (Facade)    |              |
|  +----------+                 +------+-------+              |
|                                      |                       |
|                               +------+-------+               |
|                               |DatabaseHelper|               |
|                               |   (sqflite)  |               |
|                               +--------------+               |
|                                                               |
|  +----------+    Platform     +-------------------+          |
|  | Android  | <------------  | FocusShieldSvc    |          |
|  | Native   |                | (DND Foreground)  |          |
|  | (Kotlin) | ------------>  |                   |          |
|  +----------+  MethodChannel +-------------------+          |
|                                                               |
|  +-------------------+    15-min     +-------------------+     |
|  | Workmanager Isolate| <----------  | SyncWorkerService |     |
|  |  (Background)      |              |  (Orchestrator)   |     |
|  +--------+-----------+             +---------+---------+     |
|           |                                  |                 |
|           |         +-----------------------+                 |
|           |         |      SyncService (singleton)            |
|           |         | enqueueLocalChange / retryPending        |
|           |         +-------------------+--------------------+
|           |                             |
|           |                             v
|           |                    +-----------------+
|           |                    |   Firestore     |
|           |                    |  sync_outbox    |
|           |                    |  sync_tombstones|
|           |                    |  sync_conflicts |
|           |                    +--------+--------+
|           |                             |
|           |                             v
|           |                    +-----------------+
|           |                    | Cloud Functions |
|           |                    | syncWorker      |
|           |                    | pruneTombstones |
|           |                    | aiProxy         |
|           |                    | ocrProcess      |
|           |                    +-----------------+
|           |                                       |
|  +--------+--------+                              |
|  | OfflineQueueService| - SharedPreferences (JSON)|
|  |  (max 3 retries)   |                            |
|  +--------------------+                            |
|                                                     v
|                                          +----------------+
|                                          | Firebase Storage|
|                                          | ai-uploads/    |
|                                          | ocr-output/    |
|                                          +----------------+
+-------------------------------------------------------------+
\end{blockverbatim}

% ═══════════════════════════════════════════════════════════════════════════════
\chapter*{Appendix B: SM-2 Algorithm Pseudocode}
% ═══════════════════════════════════════════════════════════════════════════════

\begin{verbatim}
FUNCTION _sm2Schedule(quality, currentEaseFactor, currentInterval, currentRepetitions):
  newEaseFactor = max(1.3, currentEaseFactor + (0.1 - (5 - quality) * (0.08 + (5 - quality) * 0.02)))
  IF quality >= 3:
    newRepetitions = currentRepetitions + 1
  ELSE:
    newRepetitions = 0
  IF newRepetitions == 0: newInterval = 1
  ELSE IF newRepetitions == 1: newInterval = 6
  ELSE IF newRepetitions == 2: newInterval = 6
  ELSE: newInterval = clamp(round(currentInterval * newEaseFactor), 1, 365)
  dueAt = now + newInterval days
  RETURN { easeFactor: newEaseFactor, intervalDays: newInterval,
           repetitions: newRepetitions, dueAt, lastReviewAt: now }
\end{verbatim}

% ═══════════════════════════════════════════════════════════════════════════════
\chapter*{Appendix C: Firestore Security Rule Pseudocode}
% ═══════════════════════════════════════════════════════════════════════════════

\begin{verbatim}
FUNCTION isOwner(userId):
  RETURN request.auth != null AND request.auth.uid == userId

MATCH /sync_outbox/{outboxId}:
  ALLOW read, write: IF isAuthenticated() AND resource.data.userId == request.auth.uid
  ALLOW create: IF isAuthenticated() AND request.resource.data.userId == request.auth.uid

MATCH /sync_tombstones/{tombstoneId}:
  ALLOW read, write: IF isAuthenticated() AND resource.data.userId == request.auth.uid
  ALLOW create: IF isAuthenticated() AND request.resource.data.userId == request.auth.uid

MATCH /quizzes/{quizId}:
  ALLOW read: IF isAuthenticated()
  ALLOW create: IF isAuthenticated() AND request.resource.data.ownerId == request.auth.uid
  ALLOW update, delete: IF isAuthenticated() AND resource.data.ownerId == request.auth.uid
\end{verbatim}

\vfill
\begin{center}
\rule{0.5\textwidth}{0.4pt}\\[1em]
\textbf{Retainly v1.0} \textbar\ \textbf{CodeSym} \textbar\ \textbf{Apache 2.0}
\end{center}

\end{document}
